%------------------------------------------------------------------------------%
\documentclass[fleqn,utf8,aspectratio=169,12pt]{beamer}

\usepackage{calculo_varias_variaveis-1}

\title{Vetor Tangente}

%------------------------------------------------------------------------------%
\begin{document}

\addtitle

%------------------------------------------------------------------------------%
\section{Vetor Tangente}

%------------------------------------------------------------------------------%
\begin{frame}{Curva Paramétrica}
  $x$ e $y$ são dados como funções
  \[
    x = f(t)     \qquad
    y = g(t)     \qquad
    t \in I
  \]

  \pause
  Vetorialmente
  \[
    \gamma(t)
    \;=\; \vetor{x}{y}
    \;=\; \vetor{f(t)}{g(t)}
  \]
\end{frame}

\framedgraphic{Vetor Tangente}{B_04_vetor_tangente_1}
\framedgraphic{Vetor Tangente}{B_04_vetor_tangente_2}
\framedgraphic{Vetor Tangente}{B_04_vetor_tangente_3}
\framedgraphic{Vetor Tangente}{B_04_vetor_tangente_4}
\framedgraphic{Vetor Tangente}{B_04_vetor_tangente_5}
\framedgraphic{Vetor Tangente}{B_04_vetor_tangente_6}
\framedgraphic{Vetor Tangente}{B_04_vetor_tangente_7}

%------------------------------------------------------------------------------%
\begin{frame}{Vetor Tangente}
  Se
  \[
    \gamma(t)
    = \vetor{x}{y}
    = \vetor{f(t)}{g(t)}
    \qquad
    t \in I
  \]
  \pause
  e $f$ e $g$ forem deriváveis em $t$,
  \pause
  o \emph{vetor tangente} é
  \pause
  \[
    \gamma'(t)
    \pause
    \;=\; \frac{d}{dt} \vetor{x}{y}
    \pause
    \;=\; \vetor{x}{y}'
    \pause
    \;=\; \Vetor{\dfrac{df}{dt}}{\dfrac{dg}{dt}}
    \pause
    \;=\; \vetor{f'(t)}{g'(t)}
  \]
\end{frame}

%------------------------------------------------------------------------------%
\section{Exemplos}

\include{ex/B-04-Vetor_tangente-ex1}
%------------------------------------------------------------------------------%
\begin{question}
\begin{columns}
\column{0.5\linewidth}
  Use derivação implícita para
  calcular o vetor tangente da
  curva paramétrica
  \[
    \gamma(t) = \big(x(t), y(t)\big)
  \]
  definida pelas equações
  \[
    x^3 + 2t^2 = 9
    \qquad
    2y^3 - 3t^2 = 4
  \]

\column{0.5\linewidth}
  \pause
  ~\\[10mm]
  \includegraphics{B_04_exemplo_2_curva}
\end{columns}
\end{question}

%------------------------------------------------------------------------------%
\begin{resolution}{Derivação Explicita ou Implícita}
  Nesse caso podemos isolar $x$ e $y$
  \[
    x = \sqrt[3]{9-2t^2}
    \qquad
    y = \sqrt[3]{\frac{4+3t^2}{2}}
  \]

  \pause
  Como resolver se não fosse possível?
\end{resolution}

%------------------------------------------------------------------------------%
\begin{resolution}{Solução}
\begin{columns}
\column{0.5\linewidth}
	\begin{align*}
		\onslide<+->{x^3 + 2t^2 & = 9 \\}
		\onslide<+->{x^3 & = 9 - 2t^2  \\}
		\onslide<+->{\frac{d}{dt}\left(x^3\right) & = \frac{d}{dt}\left(9 - 2t^2\right) \\}
		\onslide<+->{3x^2\frac{dx}{dt} & = 0 - 2\times 2t \\}
		\onslide<+->{\frac{dx}{dt} & = \frac{-4t}{3x^2}}
	\end{align*}
\column{0.5\linewidth}
	\begin{align*}
		\onslide<+->{2y^3 - 3t^2 & = 4 \\}
		\onslide<+->{2y^3 & = 4 + 3t^2 \\}
		\onslide<+->{\frac{d}{dt}\left(2y^3\right) & = \frac{d}{dt}\left(4 + 3t^2\right) \\}
		\onslide<+->{2\times 3y^2\frac{dy}{dt} & = 0 + 3 \times 2 t \\}
		\onslide<+->{\frac{dy}{dt} & = \frac{6t}{6y^2}}
		\onslide<+->{= \frac{t}{y^2}}
	\end{align*}
	\end{columns}
\end{resolution}

%------------------------------------------------------------------------------%
\begin{resolution}{Solução}
  Vetor tangente
  \[
    \gamma'(t)
    \;=\; \Vetor{\frac{-4t}{3x^2}}{\frac{t}{y^2}}
  \]
\end{resolution}

%------------------------------------------------------------------------------%

\include{ex/B-04-Vetor_tangente-ex3}
\include{ex/B-04-Vetor_tangente-ex4}
%------------------------------------------------------------------------------%
\begin{question}
\begin{columns}
\column{0.5\linewidth}
  Encontre o vetor tangente ao astroide
  \[
    x = \cos^3(t) \qquad
  \]

  \[
    y = \sin^3(t) \qquad
  \]

  \[
    0 \leq t \leq 2\pi
  \]

  no ponto \(t=\frac{\pi}{6}\)
\column{0.5\linewidth}
  \pause
  ~\\[5mm]
  \includegraphics{B_04_exemplo_5_curva}
\end{columns}
\end{question}

%------------------------------------------------------------------------------%
\begin{resolution}{Solução}
  Avaliando das derivadas das componentes

  \pause
  \medskip
	\[
		\frac{dx}{dt}
		\pause
		\;=\; \frac{d}{dt}\cos^3(t)
		\pause
		\;=\; 3\cos^2(t)\frac{d}{dt}\cos(t)
		\pause
		\;=\; -3\cos^2(t)\sin(t)
	\]

	\pause
  \medskip
	\[
		\frac{dy}{dt}
		\pause
		\;=\; \frac{d}{dt}\sin^3(t)
		\pause
		\;=\; 3\sin^2(t)\frac{d}{dt}\sin(t)
		\pause
		\;=\; 3\sin^2(t)\cos(t)
	\]
\end{question}

%------------------------------------------------------------------------------%
\begin{resolution}{Solução}
  Avaliando a derivada de $x$ no ponto \(t=\frac{\pi}{6}\)
  \pause
  \[
    \frac{dx}{dt}\left(\frac{\pi}{6}\right)
    \pause
    \;=\; \left[ -3\cos^2(t)\sin(t)\right] \evalat{\frac{\pi}{6}}{}
    \pause
    \;=\; -3\cos^2\left(\frac{\pi}{6}\right)\sin\left(\frac{\pi}{6}\right)
  \]
  \[
    \pause
    \hphantom{\frac{dx}{dt}\left(\frac{\pi}{6}\right)}
    \;=\; -3\left(\frac{\sqrt{3}}{2}\right)^2 \left(\frac{1}{2}\right)
    \pause
    \;=\; -3\; \frac{3}{4}\; \frac{1}{2}
    \pause
    \;=\; \frac{-9}{8}
  \]
\end{resolution}

%------------------------------------------------------------------------------%
\begin{resolution}{Solução}
  Avaliando a derivada de $y$ no ponto \(t=\frac{\pi}{6}\)
  \pause
  \[
    \frac{dy}{dt}\left(\frac{\pi}{6}\right)
    \pause
    \;=\; \left[ 3\sin^2(t)\cos(t)\right] \evalat{\frac{\pi}{6}}{}
    \pause
    \;=\; 3\sin^2\left(\frac{\pi}{6}\right)\cos\left(\frac{\pi}{6}\right)
  \]
  \[
    \pause
    \hphantom{\frac{dy}{dt}\left(\frac{\pi}{6}\right)}
    \;=\; 3\left(\frac{1}{2}\right)^2 \left(\frac{\sqrt{3}}{2}\right)
    \pause
    \;=\; 3 \; \frac{1}{4} \; \frac{\sqrt{3}}{2}
    \pause
    \;=\; \frac{3\sqrt{3}}{8}
  \]
\end{resolution}

%------------------------------------------------------------------------------%
\begin{resolution}{Solução}
  Função vetor tangente
  \pause
  \[
    v(t)
    \pause
    \;=\; \frac{d}{dt}\vetor{x}{y}
    \pause
    \;=\; \vetor{-3\cos^2(t)\sin(t)}{\hpm 3\sin^2(t)\cos(t)}
  \]

  \pause
  No ponto \(t=\frac{\pi}{6}\)
  \[
    \pause
    v\left(\frac{\pi}{6}\right)
    \pause
    \;=\; \Vetor{\frac{-9}{8}}{\frac{3\sqrt{3}}{8}}
    \pause
    \;=\; \frac{3}{8} \vetor{-3}{\sqrt{3}}
  \]
\end{resolution}

%------------------------------------------------------------------------------%
\begin{resolution}{Solução}
  \centering
  \includegraphics{B_04_exemplo_5_vetor}
\end{resolution}

%------------------------------------------------------------------------------%


%------------------------------------------------------------------------------%
\section{Lista Mínima}

%------------------------------------------------------------------------------%
\begin{frame}{Lista Mínima}
  Cálculo Vol. 2 do Thomas 12\textsuperscript{a} ed. -- Seção 11.2
  \begin{enumerate}
    \item Estudar todo o texto da seção
    \item \emph{Calcule o vetor tangente}
          das funções dos exercícios: 2, 8, 12, 16, 18, 20\\
          (Ignore o enunciado original dos exercícios)
  \end{enumerate}

  \vfill
  Atenção: A prova é baseada no livro, não nas apresentações
\end{frame}

%------------------------------------------------------------------------------%
\end{document}
%------------------------------------------------------------------------------%
